\begin{statement}{C.2.UN.20}{Proved}{The rank-one component fibration descends to the original support}
In the setting and with the actual component cover of C.2.UN.16,
the generic-rank-one branch gives a fibration on a resolution of
the original normalization. More precisely, for the stabilizer
$G$ of the chosen covering component $E$, there are proper
birational maps and quotient maps in a commutative diagram
\[
\begin{array}{ccc}
 \widehat E&\xrightarrow{\ q\ }&B\\
 \downarrow&&\downarrow\\
 \widehat E/G&\longrightarrow&B/G,
\end{array}                                                    \tag{C.FD.1}
\]
where $B$ is a smooth compact curve, $q$ has connected general
fibre, and $\widehat E/G$ is bimeromorphic to $S_i^\nu$.
Resolving this last quotient and taking its Stein factor gives
an actual proper fibration
\[
        Z'_i\longrightarrow B/G
        \quad\hbox{with connected rational or elliptic general fibre}.
                                                               \tag{C.FD.2}
\]
The base in (C.FD.2) is already the Stein base; in particular
there is no additional finite cover of $B/G$ implicit in this
description. For every original component with nonnegative
Kodaira dimension, its general fibre in (C.FD.2) is elliptic.
The same conclusion holds for every original component that
is not uniruled.
\end{statement}
\begin{proof}
Use a resolution equivariant for the finite cyclic group of the
normalized base change, as in C.2.UN.19. Its stabilizer $G$
preserves $E$. Recall the actual modules
\[
 M=(f'_*\mathcal T_{U'/\Delta_s})_0,
 \qquad
 L=\{v\in M\otimes_{\mathbb C\{s\}}\mathbb C(\{s\}):
                  \operatorname{ord}_E(v)\ge0\}.
                                                               \tag{C.FD.3}
\]
An element $g\in G$ covers $s\mapsto\zeta_gs$. Pushforward of
vector fields is therefore a semilinear automorphism of $M$
and of its meromorphic localization. Since $g(E)=E$, it
preserves the valuation defining $L$. It also preserves $sL$,
because $s$ changes by a nonzero constant. Reduction gives
a complex-linear representation on $L/sL$, and the restriction
map of (C.CF.6) is equivariant:
\[
 g_*\bigl(v|_E\bigr)=(g_*v)|_E.
                                                               \tag{C.FD.4}
\]
Consequently the two-dimensional space $\mathfrak a_E$ and
its rank-one meromorphic tangent distribution are preserved
by $G$.

For a basis $V_1,V_2$ write $V_2=hV_1$. The function $h$ is
nonconstant, and both fields annihilate it. If the matrix of
$g_*$ in this basis is
\[
       g_*V_1=aV_1+bV_2,\qquad
       g_*V_2=cV_1+dV_2,
\]
then the actual coefficient comparison gives
\[
              h\circ g^{-1}=\frac{c+dh}{a+bh}.
                                                               \tag{C.FD.5}
\]
This identity holds where the displayed denominator is nonzero
and hence as an identity of meromorphic functions. It supplies
the finite projective action on the target of
$h:E\dashrightarrow\mathbb P^1$. Resolve its graph
equivariantly, and then take the Stein factor. Uniqueness of
the Stein factor lifts the action on $\mathbb P^1$ to an action
on the smooth compact curve $B$, and gives an equivariant map
$q:\widehat E\to B$. This is the same connected-leaf
fibration as in C.2.UN.16. Its general fibre is rational or
elliptic, by the semicomplete one-dimensional field calculation
there. Taking the finite quotient produces (C.FD.1).

We spell out the general quotient fibre. Let
$K\subset G$ be the kernel of the action on $B$. Outside
the finitely many fixed points of the nonidentity elements
of $G/K$, a point $b\in B$ has stabilizer exactly $K$.
The inverse image of its orbit in $B/G$ is the disjoint union
of the fibres $q^{-1}(gb)$, and $G$ acts transitively on these
fibres. Hence the quotient fibre is exactly
\[
                       q^{-1}(b)/K.
                                                               \tag{C.FD.6}
\]
It is connected. A finite quotient of a smooth compact curve
is smooth, as follows locally from the invariant coordinate
$z\mapsto z^r$ at a cyclic stabilizer. Riemann--Hurwitz
shows that the genus of (C.FD.6) is at most the genus of
$q^{-1}(b)$, and is therefore zero or one.

For clarity, the generic fibre remains unchanged when the
surface quotient is resolved. A stabilizer of a point in a
general fibre is contained in $K$. It fixes the base coordinate
and, after finite holomorphic linearization, acts by
$(x,y)\mapsto(x,\xi y)$. Its invariant ring is
$\mathbb C\{x,y^r\}$, so the surface quotient is smooth
along a general fibre. A resolution can be an isomorphism
there. The resulting map to $B/G$ has connected general
fibre by (C.FD.6). Its Stein factor is a finite map of normal
compact curves of degree one, hence an isomorphism. This
proves the assertion about the base in (C.FD.2).

The quotient comparisons (C.HQ.2)--(C.HQ.3), or directly
the generic $G$-torsor in (C.HQ.1), identify
$\widehat E/G$ bimeromorphically with $S_i^\nu$.
Resolving their graph gives the proper birational comparison
of (C.FD.2) with a smooth normalization resolution of the
original component. Thus the fibration has been constructed
on the original birational component, through the specified
finite quotient, rather than inferred from images of a
covering family of curves.

Finally a rational general fibre $F$ of (C.FD.2) has trivial
normal line in $Z'_i$, and adjunction gives
$K_{Z'_i}|_F=K_F=\mathcal O_{\mathbb P^1}(-2)$.
Every positive pluricanonical section restricts to zero on
every general fibre and therefore vanishes identically.
This makes $\kappa(Z'_i)=-\infty$, so the nonnegative
Kodaira-dimension cases have elliptic general fibre.
More generally the rational fibration itself is a family
of rational curves covering the surface; it is precisely
an actual ruling on a proper birational model and is
excluded when the original component is not uniruled.
\end{proof}


% C.2.UN.26 from part_C2_field_limit.tex

